%!TEX program = xelatex
\documentclass[11pt,a4paper]{article}
\usepackage[margin=2.2cm]{geometry}
\usepackage{amsmath,amssymb,amsthm,mathtools}
\usepackage{bm}
\usepackage{enumitem}
\usepackage{booktabs}
\usepackage{xcolor}
\usepackage{hyperref}
\usepackage{fancyhdr}
\usepackage{xeCJK}

\hypersetup{
    colorlinks=true,
    linkcolor=blue!60!black,
    citecolor=blue!60!black,
    urlcolor=blue!60!black
}

\pagestyle{fancy}
\fancyhf{}
\rhead{\small Qiuzhen QE Computational \& Applied Math --- Fall 2023}
\lhead{\small Official Qualifying Exam}
\rfoot{\small Page \thepage}

\DeclareMathOperator*{\argmax}{arg\,max}
\DeclareMathOperator*{\argmin}{arg\,min}

\setlist[itemize]{topsep=3pt,itemsep=2pt,parsep=0pt}
\setlist[enumerate]{topsep=3pt,itemsep=2pt,parsep=0pt}

\newcommand{\examstatement}[1]{%
  \par\addvspace{0.85em}%
  \noindent{\bfseries #1}\par\nobreak\addvspace{0.28em}%
}

% Shared amsthm note/remark/theorem styles
% Shared math extras for qe-review.
% Requires already-loaded: amsmath, amssymb.
%
% Classic pitfall: AMS blackboard-bold fonts have no digit glyphs, so
% \mathbb{1} renders as overlapping garbage. Map the digit 1 to dsfont's
% \mathds{1}; leave \mathbb{R}, \mathbb{P}, etc. on AMS.
%
% Usage (path relative to the main .tex being compiled):
%   notes:  % Shared math extras for qe-review.
% Requires already-loaded: amsmath, amssymb.
%
% Classic pitfall: AMS blackboard-bold fonts have no digit glyphs, so
% \mathbb{1} renders as overlapping garbage. Map the digit 1 to dsfont's
% \mathds{1}; leave \mathbb{R}, \mathbb{P}, etc. on AMS.
%
% Usage (path relative to the main .tex being compiled):
%   notes:  % Shared math extras for qe-review.
% Requires already-loaded: amsmath, amssymb.
%
% Classic pitfall: AMS blackboard-bold fonts have no digit glyphs, so
% \mathbb{1} renders as overlapping garbage. Map the digit 1 to dsfont's
% \mathds{1}; leave \mathbb{R}, \mathbb{P}, etc. on AMS.
%
% Usage (path relative to the main .tex being compiled):
%   notes:  \input{../../common/qe-math-extras.tex}
%   exams:  \input{../../../common/qe-math-extras.tex}

\usepackage{dsfont}
\usepackage{pdftexcmds}

\makeatletter
\let\qe@amsmmathbb\mathbb
\renewcommand{\mathbb}[1]{%
  \ifnum\pdf@strcmp{\unexpanded{#1}}{1}=\z@
    \mathds{1}%
  \else
    \qe@amsmmathbb{#1}%
  \fi
}
\makeatother

% Preferred indicator macros (either style is safe after the remap above).
\newcommand{\bbone}{\mathds{1}}
\newcommand{\ind}{\mathbf{1}}

%   exams:  % Shared math extras for qe-review.
% Requires already-loaded: amsmath, amssymb.
%
% Classic pitfall: AMS blackboard-bold fonts have no digit glyphs, so
% \mathbb{1} renders as overlapping garbage. Map the digit 1 to dsfont's
% \mathds{1}; leave \mathbb{R}, \mathbb{P}, etc. on AMS.
%
% Usage (path relative to the main .tex being compiled):
%   notes:  \input{../../common/qe-math-extras.tex}
%   exams:  \input{../../../common/qe-math-extras.tex}

\usepackage{dsfont}
\usepackage{pdftexcmds}

\makeatletter
\let\qe@amsmmathbb\mathbb
\renewcommand{\mathbb}[1]{%
  \ifnum\pdf@strcmp{\unexpanded{#1}}{1}=\z@
    \mathds{1}%
  \else
    \qe@amsmmathbb{#1}%
  \fi
}
\makeatother

% Preferred indicator macros (either style is safe after the remap above).
\newcommand{\bbone}{\mathds{1}}
\newcommand{\ind}{\mathbf{1}}


\usepackage{dsfont}
\usepackage{pdftexcmds}

\makeatletter
\let\qe@amsmmathbb\mathbb
\renewcommand{\mathbb}[1]{%
  \ifnum\pdf@strcmp{\unexpanded{#1}}{1}=\z@
    \mathds{1}%
  \else
    \qe@amsmmathbb{#1}%
  \fi
}
\makeatother

% Preferred indicator macros (either style is safe after the remap above).
\newcommand{\bbone}{\mathds{1}}
\newcommand{\ind}{\mathbf{1}}

%   exams:  % Shared math extras for qe-review.
% Requires already-loaded: amsmath, amssymb.
%
% Classic pitfall: AMS blackboard-bold fonts have no digit glyphs, so
% \mathbb{1} renders as overlapping garbage. Map the digit 1 to dsfont's
% \mathds{1}; leave \mathbb{R}, \mathbb{P}, etc. on AMS.
%
% Usage (path relative to the main .tex being compiled):
%   notes:  % Shared math extras for qe-review.
% Requires already-loaded: amsmath, amssymb.
%
% Classic pitfall: AMS blackboard-bold fonts have no digit glyphs, so
% \mathbb{1} renders as overlapping garbage. Map the digit 1 to dsfont's
% \mathds{1}; leave \mathbb{R}, \mathbb{P}, etc. on AMS.
%
% Usage (path relative to the main .tex being compiled):
%   notes:  \input{../../common/qe-math-extras.tex}
%   exams:  \input{../../../common/qe-math-extras.tex}

\usepackage{dsfont}
\usepackage{pdftexcmds}

\makeatletter
\let\qe@amsmmathbb\mathbb
\renewcommand{\mathbb}[1]{%
  \ifnum\pdf@strcmp{\unexpanded{#1}}{1}=\z@
    \mathds{1}%
  \else
    \qe@amsmmathbb{#1}%
  \fi
}
\makeatother

% Preferred indicator macros (either style is safe after the remap above).
\newcommand{\bbone}{\mathds{1}}
\newcommand{\ind}{\mathbf{1}}

%   exams:  % Shared math extras for qe-review.
% Requires already-loaded: amsmath, amssymb.
%
% Classic pitfall: AMS blackboard-bold fonts have no digit glyphs, so
% \mathbb{1} renders as overlapping garbage. Map the digit 1 to dsfont's
% \mathds{1}; leave \mathbb{R}, \mathbb{P}, etc. on AMS.
%
% Usage (path relative to the main .tex being compiled):
%   notes:  \input{../../common/qe-math-extras.tex}
%   exams:  \input{../../../common/qe-math-extras.tex}

\usepackage{dsfont}
\usepackage{pdftexcmds}

\makeatletter
\let\qe@amsmmathbb\mathbb
\renewcommand{\mathbb}[1]{%
  \ifnum\pdf@strcmp{\unexpanded{#1}}{1}=\z@
    \mathds{1}%
  \else
    \qe@amsmmathbb{#1}%
  \fi
}
\makeatother

% Preferred indicator macros (either style is safe after the remap above).
\newcommand{\bbone}{\mathds{1}}
\newcommand{\ind}{\mathbf{1}}


\usepackage{dsfont}
\usepackage{pdftexcmds}

\makeatletter
\let\qe@amsmmathbb\mathbb
\renewcommand{\mathbb}[1]{%
  \ifnum\pdf@strcmp{\unexpanded{#1}}{1}=\z@
    \mathds{1}%
  \else
    \qe@amsmmathbb{#1}%
  \fi
}
\makeatother

% Preferred indicator macros (either style is safe after the remap above).
\newcommand{\bbone}{\mathds{1}}
\newcommand{\ind}{\mathbf{1}}


\usepackage{dsfont}
\usepackage{pdftexcmds}

\makeatletter
\let\qe@amsmmathbb\mathbb
\renewcommand{\mathbb}[1]{%
  \ifnum\pdf@strcmp{\unexpanded{#1}}{1}=\z@
    \mathds{1}%
  \else
    \qe@amsmmathbb{#1}%
  \fi
}
\makeatother

% Preferred indicator macros (either style is safe after the remap above).
\newcommand{\bbone}{\mathds{1}}
\newcommand{\ind}{\mathbf{1}}

% Shared amsthm environments for qe-review (minimal).
% Requires already-loaded: amsthm.
% Safe to \input after \usepackage{amsmath,amssymb,amsthm,...}.
%
% Shared numbering uses aliascnt so cleveref keys off the environment
% name (Lemma, Proposition, Exercise, ...) rather than the parent
% counter (Theorem, Definition, Remark). Without aliascnt, \cref{lem:...}
% prints ``Theorem 9.13'' instead of ``Lemma 9.13''.

\usepackage{aliascnt}

\theoremstyle{plain}
\newtheorem{theorem}{Theorem}[section]

\newaliascnt{lemma}{theorem}
\newtheorem{lemma}[lemma]{Lemma}
\aliascntresetthe{lemma}

\newaliascnt{proposition}{theorem}
\newtheorem{proposition}[proposition]{Proposition}
\aliascntresetthe{proposition}

\newaliascnt{corollary}{theorem}
\newtheorem{corollary}[corollary]{Corollary}
\aliascntresetthe{corollary}

\newaliascnt{claim}{theorem}
\newtheorem{claim}[claim]{Claim}
\aliascntresetthe{claim}

\newaliascnt{fact}{theorem}
\newtheorem{fact}[fact]{Fact}
\aliascntresetthe{fact}

\theoremstyle{definition}
\newtheorem{definition}{Definition}[section]

\newaliascnt{exercise}{definition}
\newtheorem{exercise}[exercise]{Exercise}
\aliascntresetthe{exercise}

\newaliascnt{assumption}{definition}
\newtheorem{assumption}[assumption]{Assumption}
\aliascntresetthe{assumption}

\newaliascnt{notation}{definition}
\newtheorem{notation}[notation]{Notation}
\aliascntresetthe{notation}

\newaliascnt{construction}{definition}
\newtheorem{construction}[construction]{Construction}
\aliascntresetthe{construction}

% Example: document-wide numbering (not per section / chapter)
\newtheorem{example}{Example}

\theoremstyle{remark}
\newtheorem{remark}{Remark}[section]
\newtheorem*{remark*}{Remark}

\newaliascnt{heuristic}{remark}
\newtheorem{heuristic}[heuristic]{Heuristic}
\aliascntresetthe{heuristic}

\newaliascnt{convention}{remark}
\newtheorem{convention}[convention]{Convention}
\aliascntresetthe{convention}

\newaliascnt{warning}{remark}
\newtheorem{warning}[warning]{Warning}
\aliascntresetthe{warning}

\newaliascnt{proofsketch}{remark}
\newtheorem{proofsketch}[proofsketch]{Proof sketch}
\aliascntresetthe{proofsketch}

\newaliascnt{checkpoint}{remark}
\newtheorem{checkpoint}[checkpoint]{Checkpoint}
\aliascntresetthe{checkpoint}

% Note: document-wide numbering (not per section)
\newtheorem{note}{Note}
\newtheorem*{note*}{Note}

\newaliascnt{examnote}{note}
\newtheorem{examnote}[examnote]{Exam note}
\aliascntresetthe{examnote}

\newaliascnt{study}{note}
\newtheorem{study}[study]{Study note}
\aliascntresetthe{study}

\newaliascnt{summary}{note}
\newtheorem{summary}[summary]{Summary}
\aliascntresetthe{summary}

\newtheorem{examproblem}{Problem}[section]
\newtheorem*{examproblem*}{Problem}

% Bilingual aliases
\newenvironment{dingli}{\begin{theorem}}{\end{theorem}}
\newenvironment{yinli}{\begin{lemma}}{\end{lemma}}
\newenvironment{mingti}{\begin{proposition}}{\end{proposition}}
\newenvironment{tuilun}{\begin{corollary}}{\end{corollary}}
\newenvironment{dingyi}{\begin{definition}}{\end{definition}}
\newenvironment{li}{\begin{example}}{\end{example}}
\newenvironment{zhu}{\begin{remark}}{\end{remark}}
\newenvironment{biji}{\begin{note}}{\end{note}}
\newenvironment{kaoshi}{\begin{examnote}}{\end{examnote}}

\renewcommand{\qedsymbol}{\ensuremath{\square}}

% PDF sidebar outline + printed TOC for QE exam transcripts.
% Load in the preamble AFTER hyperref and AFTER \newcommand{\examstatement}.
\hypersetup{
  unicode=true,
  bookmarks=true,
  bookmarksopen=true,
  bookmarksopenlevel=3,
  bookmarksnumbered=false,
  bookmarksdepth=3
}
\IfFileExists{bookmark.sty}{%
  \usepackage{bookmark}%
  \bookmarksetup{open,openlevel=3,numbered=false,depth=3}%
}{}
% Unnumbered in print, but still written to the TOC / PDF outline
% down to \subsubsection. Starred headings create neither.
\setcounter{secnumdepth}{-1}
\setcounter{tocdepth}{3}
\makeatletter
\@ifundefined{examstatement}{}{%
  \let\qe@origexamstatement\examstatement
  \renewcommand{\examstatement}[1]{%
    \par\phantomsection
    \addcontentsline{toc}{subsection}{#1}%
    \qe@origexamstatement{#1}%
  }%
}
\makeatother


\title{\textbf{QUALIFYING EXAM: APPLIED MATHEMATICS}\\[0.5em]\Large SEPTEMBER 2023}
\date{}
\author{}

\begin{document}

\maketitle
\vspace{-3em}

\tableofcontents

\section{Statement of the problems}

\examstatement{Problem 1.}
Consider the Newton's method for finding a solution \(x_{*}\) to \(f(x) = 0\), where \(f \in C^{2}(a, b)\), \(x_{*} \in (a, b)\).
\begin{enumerate}[label=\arabic*:,leftmargin=2em]
    \item determine \(x_{0} \in (a, b)\).
    \item for \(k = 0, 1, 2, \dots\) do
    \item \(x_{k+1} = x_{k} - \dfrac{f(x_{k})}{f'(x_{k})}\).
    \item end for
\end{enumerate}
\begin{enumerate}[label=(\arabic*)]
    \item Prove that if \(x_{0}\) is sufficiently close to \(x_{*}\) and \(f'(x_{*}) \neq 0\), then \(\lim_{k \to \infty} x_{k} = x_{*}\) and
    \[
    \lim_{k \to \infty} \frac{x_{k+1} - x_{*}}{(x_{k} - x_{*})^{2}} = \frac{f''(x_{*})}{2 f'(x_{*})}.
    \]
    \item In practice sometimes the derivative is not easy to be obtained. As a result, a difference is used instead:
    \[
    x_{k+1} = x_{k} - \frac{x_{k} - x_{k-1}}{f(x_{k}) - f(x_{k-1})} f(x_{k}).
    \]
    Prove that if \(x_{0}\) is sufficiently close to \(x_{*}\), then \(x_{k} \to x_{*}\) and
    \[
    \lim_{k \to \infty} \frac{x_{k+1} - x_{*}}{(x_{k} - x_{*})(x_{k-1} - x_{*})} = \frac{f''(x_{*})}{2 f'(x_{*})}.
    \]
\end{enumerate}

\examstatement{Problem 2.}
Consider the explicit shifted QR method for computing eigenvalues of a matrix \(A \in \mathbb{C}^{n \times n}\).
\begin{enumerate}[label=\arabic*:,leftmargin=2em]
    \item find an upper Hessenberg matrix \(H_{0}\) and a unitary matrix \(U_{0}\) such that \(H_{0} = U_{0}^{H} A U_{0}\).
    \item for \(i = 0, 1, 2, \dots\) do
    \item determine a scalar \(\mu_{k}\).
    \item compute QR factorization \(Q_{k} R_{k} = H_{k} - \mu_{k} I\).
    \item \(H_{k+1} = R_{k} Q_{k} + \mu_{k} I\).
    \item end for
\end{enumerate}
\begin{enumerate}[label=(\arabic*)]
    \item Prove that \(H_{i}\), \(i = 1, 2, \dots\) are all upper Hessenberg matrices.
    \item Interpret the purpose to use \(H_{0}\) rather than \(A\) in the iteration.
    \item Suppose that \(A\) has \(n\) distinct eigenvalues and none of the shifts \(\mu_{i}\), \(i = 1, 2, \dots\) is an eigenvalue of \(A\). Prove that \(H_{i}\), \(i = 0, 1, 2, \dots\) are unreduced upper Hessenberg matrices. (An upper Hessenberg matrix \(H\) is called unreduced, if \(H_{i+1,i} \neq 0\) for \(i = 1, \dots, n-1\).)
    \item Write \(H_{k} = \begin{bmatrix} G_{k} & u_{k} \\ \varepsilon_{k} e^{T} & \alpha_{k} \end{bmatrix}\) where \(\alpha_{k}, \varepsilon_{k} \in \mathbb{C}\), \(u_{k}, e \in \mathbb{C}^{n-1}\) and \(e = [0, \dots, 0, 1]^{T}\). Suppose \(A\) has \(n\) distinct eigenvalues. Prove that
    \[
    |\varepsilon_{k+1}| \le \rho_{k}^{2} \|u_{k}\|_{2} |\varepsilon_{k}|^{2} + \rho_{k} |\alpha_{k} - \mu_{k}| |\varepsilon_{k}|,
    \]
    where \(\rho_{k} = \|(G_{k} - \mu_{k} I)^{-1}\|_{2}\), provided \(\mu_{k}\) is not an eigenvalue of \(G_{k}\).
\end{enumerate}

\examstatement{Problem 3.}
If \(p\) is a polynomial of degree \(n\) (on \([-1, 1]\)), it is determined by its values on an \((n+1)\)-point grid on \([-1, 1]\). The derivative \(p'\), a polynomial of degree \((n-1)\), is determined on the same grid. The (classical) differentiation matrix is the \((n+1)\)-by-\((n+1)\) matrix \(D = (D_{ij}) \in \mathbb{R}^{(n+1) \times (n+1)}\) that represents the linear map from the vector of values of \(p\) to the vector of values of \(p'\), namely:

\[
p'(x_{i}) = \sum_{j=0}^{n} D_{ij} p(x_{j}).
\]

\begin{enumerate}[label=(\arabic*)]
    \item Prove that \(D_{ij} = l_{j}'(x_{i})\), where \(l_{j}(x)\) is the \(j\)-th Lagrange basis function.
    \item If a Chebyshev grid (\(x_{j} = \cos(j\pi/n)\), \(0 \le j \le n\)) is adopted, derive explicit formulas for \(D_{ij}\).
\end{enumerate}

\examstatement{Problem 4.}
Consider a Runge--Kutta method for the ODE \(y' = f(t, y)\) (\(f\) is Lipschitz continuous in \(y\) and uniform in \(t\)) with the following Butcher tableau:

\[
\begin{array}{c|cc}
0 & 0 \\
1 & 1 & 0 \\
\hline
& 1/2 & 1/2
\end{array}
\]

\begin{enumerate}[label=(\arabic*)]
    \item Rewrite the scheme in the form of \(u_{n+1} = u_{n} + h F(t_{n}, u_{n}, h; f)\).
    \item Assume that \(f\) is sufficiently smooth. Prove that this method is convergent and determine the order of convergence.
    \item What is the region of absolute stability of this method? Give a description that is as explicit as possible.
\end{enumerate}

\examstatement{Problem 5.}
For the equation \(u_{t} + a u_{xxx} = 0\) (\(a\) is a constant), applying the idea of the Lax--Friedrichs scheme, one can get the scheme

\[
u_{m}^{n+1} = \tfrac12 (u_{m+1}^{n} + u_{m-1}^{n}) - \tfrac12 a k h^{-3} (u_{m+2}^{n} - 2u_{m+1}^{n} + 2u_{m-1}^{n} - u_{m-2}^{n}),
\]

where \(k\) and \(h\) represent the time step and mesh size, respectively.
\begin{enumerate}[label=(\arabic*)]
    \item Give the leading order term of local truncation error.
    \item Analyze the stability of this scheme.
\end{enumerate}

\examstatement{Problem 6.}
Consider the following linear programming problem

\[
\max_{x} \langle c, x \rangle \quad \text{s.t.} \quad Ax = b,\ x \ge 0,
\]

where \(x \in \mathbb{R}^{n}\), \(0 \le b \in \mathbb{R}^{m}\), \(A \in \mathbb{R}^{m \times n}\) and \(c \in \mathbb{R}^{n}\), \(A\) is a full rank matrix and \(m < n\).
\begin{enumerate}[label=(\arabic*)]
    \item Prove that the basic feasible solutions are equivalent to the vertices of its feasible region.
    \item Write down the dual problem, and prove that when the optimal solution exists, the dual problem must also have an optimal solution, and the optimal objective values of these two problems are equal.
\end{enumerate}

\examstatement{Problem 7.}
Consider the eigenvalue problem with \(0 < \epsilon \ll 1\)

\[
u'' + (\lambda + \epsilon f(x)) u = 0, \qquad 0 < x < 1,
\]

\[
u(0) = 0, \qquad u'(1) = 0,
\]

where \(f\) is a given smooth function. Give the asymptotic expansion of \(\lambda\) such that the accuracy is \(O(\epsilon)\).

\noindent\rule{\linewidth}{0.4pt}

\section{Statement and solutions}

\subsection{Problem 1.}
Consider the Newton's method for finding a solution \(x_{*}\) to \(f(x) = 0\), where \(f \in C^{2}(a, b)\), \(x_{*} \in (a, b)\).
\begin{enumerate}[label=\arabic*:,leftmargin=2em]
    \item determine \(x_{0} \in (a, b)\).
    \item for \(k = 0, 1, 2, \dots\) do
    \item \(x_{k+1} = x_{k} - \dfrac{f(x_{k})}{f'(x_{k})}\).
    \item end for
\end{enumerate}
\begin{enumerate}[label=(\arabic*)]
    \item Prove that if \(x_{0}\) is sufficiently close to \(x_{*}\) and \(f'(x_{*}) \neq 0\), then \(\lim_{k \to \infty} x_{k} = x_{*}\) and
    \[
    \lim_{k \to \infty} \frac{x_{k+1} - x_{*}}{(x_{k} - x_{*})^{2}} = \frac{f''(x_{*})}{2 f'(x_{*})}.
    \]
    \item In practice sometimes the derivative is not easy to be obtained. As a result, a difference is used instead:
    \[
    x_{k+1} = x_{k} - \frac{x_{k} - x_{k-1}}{f(x_{k}) - f(x_{k-1})} f(x_{k}).
    \]
    Prove that if \(x_{0}\) is sufficiently close to \(x_{*}\), then \(x_{k} \to x_{*}\) and
    \[
    \lim_{k \to \infty} \frac{x_{k+1} - x_{*}}{(x_{k} - x_{*})(x_{k-1} - x_{*})} = \frac{f''(x_{*})}{2 f'(x_{*})}.
    \]
\end{enumerate}

\subsection{Problem 2.}
Consider the explicit shifted QR method for computing eigenvalues of a matrix \(A \in \mathbb{C}^{n \times n}\).
\begin{enumerate}[label=\arabic*:,leftmargin=2em]
    \item find an upper Hessenberg matrix \(H_{0}\) and a unitary matrix \(U_{0}\) such that \(H_{0} = U_{0}^{H} A U_{0}\).
    \item for \(i = 0, 1, 2, \dots\) do
    \item determine a scalar \(\mu_{k}\).
    \item compute QR factorization \(Q_{k} R_{k} = H_{k} - \mu_{k} I\).
    \item \(H_{k+1} = R_{k} Q_{k} + \mu_{k} I\).
    \item end for
\end{enumerate}
\begin{enumerate}[label=(\arabic*)]
    \item Prove that \(H_{i}\), \(i = 1, 2, \dots\) are all upper Hessenberg matrices.
    \item Interpret the purpose to use \(H_{0}\) rather than \(A\) in the iteration.
    \item Suppose that \(A\) has \(n\) distinct eigenvalues and none of the shifts \(\mu_{i}\), \(i = 1, 2, \dots\) is an eigenvalue of \(A\). Prove that \(H_{i}\), \(i = 0, 1, 2, \dots\) are unreduced upper Hessenberg matrices. (An upper Hessenberg matrix \(H\) is called unreduced, if \(H_{i+1,i} \neq 0\) for \(i = 1, \dots, n-1\).)
    \item Write \(H_{k} = \begin{bmatrix} G_{k} & u_{k} \\ \varepsilon_{k} e^{T} & \alpha_{k} \end{bmatrix}\) where \(\alpha_{k}, \varepsilon_{k} \in \mathbb{C}\), \(u_{k}, e \in \mathbb{C}^{n-1}\) and \(e = [0, \dots, 0, 1]^{T}\). Suppose \(A\) has \(n\) distinct eigenvalues. Prove that
    \[
    |\varepsilon_{k+1}| \le \rho_{k}^{2} \|u_{k}\|_{2} |\varepsilon_{k}|^{2} + \rho_{k} |\alpha_{k} - \mu_{k}| |\varepsilon_{k}|,
    \]
    where \(\rho_{k} = \|(G_{k} - \mu_{k} I)^{-1}\|_{2}\), provided \(\mu_{k}\) is not an eigenvalue of \(G_{k}\).
\end{enumerate}

\subsection{Problem 3.}
If \(p\) is a polynomial of degree \(n\) (on \([-1, 1]\)), it is determined by its values on an \((n+1)\)-point grid on \([-1, 1]\). The derivative \(p'\), a polynomial of degree \((n-1)\), is determined on the same grid. The (classical) differentiation matrix is the \((n+1)\)-by-\((n+1)\) matrix \(D = (D_{ij}) \in \mathbb{R}^{(n+1) \times (n+1)}\) that represents the linear map from the vector of values of \(p\) to the vector of values of \(p'\), namely:

\[
p'(x_{i}) = \sum_{j=0}^{n} D_{ij} p(x_{j}).
\]

\begin{enumerate}[label=(\arabic*)]
    \item Prove that \(D_{ij} = l_{j}'(x_{i})\), where \(l_{j}(x)\) is the \(j\)-th Lagrange basis function.
    \item If a Chebyshev grid (\(x_{j} = \cos(j\pi/n)\), \(0 \le j \le n\)) is adopted, derive explicit formulas for \(D_{ij}\).
\end{enumerate}

\subsection{Problem 4.}
Consider a Runge--Kutta method for the ODE \(y' = f(t, y)\) (\(f\) is Lipschitz continuous in \(y\) and uniform in \(t\)) with the following Butcher tableau:

\[
\begin{array}{c|cc}
0 & 0 \\
1 & 1 & 0 \\
\hline
& 1/2 & 1/2
\end{array}
\]

\begin{enumerate}[label=(\arabic*)]
    \item Rewrite the scheme in the form of \(u_{n+1} = u_{n} + h F(t_{n}, u_{n}, h; f)\).
    \item Assume that \(f\) is sufficiently smooth. Prove that this method is convergent and determine the order of convergence.
    \item What is the region of absolute stability of this method? Give a description that is as explicit as possible.
\end{enumerate}

\subsection{Problem 5.}
For the equation \(u_{t} + a u_{xxx} = 0\) (\(a\) is a constant), applying the idea of the Lax--Friedrichs scheme, one can get the scheme

\[
u_{m}^{n+1} = \tfrac12 (u_{m+1}^{n} + u_{m-1}^{n}) - \tfrac12 a k h^{-3} (u_{m+2}^{n} - 2u_{m+1}^{n} + 2u_{m-1}^{n} - u_{m-2}^{n}),
\]

where \(k\) and \(h\) represent the time step and mesh size, respectively.
\begin{enumerate}[label=(\arabic*)]
    \item Give the leading order term of local truncation error.
    \item Analyze the stability of this scheme.
\end{enumerate}

\subsection{Problem 6.}
Consider the following linear programming problem

\[
\max_{x} \langle c, x \rangle \quad \text{s.t.} \quad Ax = b,\ x \ge 0,
\]

where \(x \in \mathbb{R}^{n}\), \(0 \le b \in \mathbb{R}^{m}\), \(A \in \mathbb{R}^{m \times n}\) and \(c \in \mathbb{R}^{n}\), \(A\) is a full rank matrix and \(m < n\).
\begin{enumerate}[label=(\arabic*)]
    \item Prove that the basic feasible solutions are equivalent to the vertices of its feasible region.
    \item Write down the dual problem, and prove that when the optimal solution exists, the dual problem must also have an optimal solution, and the optimal objective values of these two problems are equal.
\end{enumerate}

\subsection{Problem 7.}
Consider the eigenvalue problem with \(0 < \epsilon \ll 1\)

\[
u'' + (\lambda + \epsilon f(x)) u = 0, \qquad 0 < x < 1,
\]

\[
u(0) = 0, \qquad u'(1) = 0,
\]

where \(f\) is a given smooth function. Give the asymptotic expansion of \(\lambda\) such that the accuracy is \(O(\epsilon)\).

\end{document}
